\documentclass[11pt]{letter}
\usepackage[utf8]{inputenc}
\usepackage[spanish]{babel}
\usepackage[margin=1in]{geometry}
\usepackage{parskip}

\begin{document}

\begin{flushright}
Bogotá, D.C., 31 de agosto de 2026
\end{flushright}

\textbf{Señores} \\
\textbf{Dirección de Relaciones Exteriores -- DRE} \\
Universidad Nacional de Colombia

\vspace{0.5cm}

\textbf{Asunto:} Carta de motivación -- Convocatoria 2026-0054, Estudiante Auxiliar para el apoyo al desarrollo e implementación de la página web de la Red de Embajadores de la Diáspora Académica y Científica UNAL

\vspace{0.3cm}

Respetados señores,

\vspace{0.3cm}

Por medio de la presente, manifiesto mi interés en participar en la convocatoria 2026-0054, dirigida a seleccionar un(a) estudiante auxiliar que apoye el diseño, desarrollo e implementación de herramientas digitales para la estrategia de la Red Diáspora Académica-Científica de la Universidad Nacional de Colombia.

Soy estudiante de Ciencia de la Computación en la Universidad Nacional de Colombia. A lo largo de mi formación he desarrollado experiencia práctica tanto en desarrollo web como en análisis y manejo de datos, dos frentes que considero directamente relevantes para el proyecto de la convocatoria. He trabajado con Angular para construir interfaces de usuario, y también con Python, Java y JavaScript en distintos contextos; adicionalmente, he construido páginas web utilizando HTML y CSS de forma nativa, sin depender de frameworks, lo que me ha dado una comprensión sólida de los fundamentos del desarrollo front-end.

Entre los proyectos que he desarrollado se destaca una plataforma web educativa para la divulgación del teorema de Pitágoras, en la cual los usuarios pueden manipular e interactuar con figuras geométricas para comprender el teorema de forma visual e intuitiva. Para este proyecto utilicé HTML y CSS en la interfaz, y Python junto con SQLite para la lógica del backend. De forma paralela, llevo más de un año desarrollando una herramienta de trading que hace un uso intensivo del análisis de bases de datos con Python, lo cual ha fortalecido mis habilidades para trabajar con información estructurada, procesar datos de manera eficiente y construir soluciones robustas y funcionales.

Me motiva especialmente la posibilidad de aplicar los conocimientos que he adquirido en proyectos académicos y personales a una iniciativa real, con un impacto directo en la comunidad universitaria. Participar en el desarrollo de la página web de la Red de Embajadores de la Diáspora Académica y Científica UNAL representa para mí la oportunidad de contribuir a una herramienta que fortalezca los vínculos entre la Universidad y su comunidad académica y científica en el exterior, al tiempo que sigo consolidando mis competencias técnicas en un entorno de trabajo real.

Agradezco de antemano la consideración de mi postulación y quedo atento(a) a cualquier información adicional que se requiera para el proceso de selección.

\vspace{0.3cm}

Cordialmente,

\vspace{1.5cm}

\rule{7cm}{0.4pt} \\
Nombre completo \\
Estudiante de Ciencia de la Computación \\
Universidad Nacional de Colombia

\end{document}
